\section{Removing the projective ambiguity}\label{sec:projective}

\S\ref{sec:point} matches projective normal forms, and leaves standing the
qualifier ``up to a rank-$1$ twist''. This section removes it. The removal is
not by an abstract uniqueness argument --- the obvious such argument is false,
and is refuted below --- but by exhibiting the isomorphism directly.

\subsection{Trace-freeness does not do the work}\label{ssec:tracefree}

The natural attempt is to argue that van der Put--Saito's normalisation picks a
canonical representative, so that a projective match is automatically a match.
Half of that is right.

\paragraph{What is true (\texttt{S17-01}, \textbf{PROVEN}).}
A rank-$1$ twist acts by $A\mapsto A+\lambda\,\mathrm{Id}$ and changes the
trace by $2\lambda$, so in characteristic $0$ the equation
$\operatorname{tr}(A+\lambda\,\mathrm{Id})=0$ has the unique solution
$\lambda=-\operatorname{tr}(A)/2$. Within the orbit of a fixed connection under
rank-$1$ twists there is therefore exactly one trace-free member.

\paragraph{What is false (\texttt{S17-02}, \textbf{PROVEN}).}\label{par:tracefree-false}
Trace-freeness alone does \emph{not} single out a representative of a
projective class, and the obstruction is explicit. Over $\mathbb C(z)$ put
\[
  L=\partial-\frac{1}{2z}.
\]
Then $L\otimes L=\partial-1/z$ is trivial, with rational witness $w=z$
satisfying $2f+w'/w=0$; whereas $L$ itself is non-trivial, since the residue of
$f$ at $0$ is $-\tfrac12$, not an integer, so $f$ is not the logarithmic
derivative of a rational function. Tensoring a trace-free module by such a
$2$-torsion $L$ preserves both the projective class and the triviality of the
determinant. A projective class can therefore contain more than one trace-free
module up to isomorphism.

\begin{remark}[A printed line that contradicted its own conclusion;
\texttt{S17-03}, \textbf{PROVEN}]\label{rem:contradiction}
The first run of the check above printed ``$L\otimes L$ trivial? False'' while
the surrounding text drew the opposite conclusion. The test had compared $2f$
against $-v'/v$ with $v=1/u$, rather than testing whether $2f$ is a logarithmic
derivative at all; the correct witness is $w=z$, not $1/z$. It was replaced by
a residue-based test --- all poles simple, all residues integers, witness
constructed and verified --- which is self-checked in both directions, on
$\partial-1/z$ (trivial) and on $\partial-1/(2z)$ (non-trivial), before being
used. The conclusion did not change, but the evidence as first printed did not
support it, and the episode is recorded rather than tidied away.
\end{remark}

Because the general statement fails, the correct hypothesis under which a
projective class has a unique normalised representative is not settled in this
paper. It is stated as an open question in \S\ref{sec:open}
(\texttt{S17-07}, \textbf{CONJECTURED}), and nothing below depends on it.

\subsection{Comparing the pair \texorpdfstring{$(P,Q)$}{(P,Q)}, not the
invariant \texorpdfstring{$R$}{R}}\label{ssec:pq}

A second-order operator is determined by the pair $(P,Q)$ in
$\partial^2+P\partial+Q$, and the projective invariant
\[
  R=\tfrac14P^2+\tfrac12P'-Q
\]
that \S\ref{sec:point} matched is by construction blind to a rank-$1$ twist.
The pair is not. So the ambiguity is removed by comparing both coefficients
rather than their invariant combination --- which requires no new theory, only
that the comparison actually be made.

It has been (\textbf{PROVEN}, \texttt{S17-04}). In the $z$ coordinate
System~$\mathsf C$ has
\begin{align}
  P_{\mathsf C}&=\frac{2z-1}{z(z-1)},\label{eq:PC}\\[2pt]
  Q_{\mathsf C}&=\frac{22z^2-22z-2i\sqrt{11}\,z+5+i\sqrt{11}}
                     {54\,z(z-1)},\label{eq:QC}
\end{align}
and the family's scalar operator at the located point $(p,q,t,\theta_\infty)$,
taken with cyclic vector $e_1$, has \emph{exactly} the same $P$ and the same
$Q$: both differences vanish identically, with worst relative residual $0$ over
five sample points at $60$ digits. System~$\mathsf C$'s differential module and
the family's module at that point are therefore isomorphic, not merely
projectively isomorphic.

Both vacuity controls fired. A value of $q$ perturbed by $10^{-3}$ is rejected;
and an operator constructed to have the same $R$ but a different $P$ is
separated by the $(P,Q)$ test though not by $R$ --- which is the precise sense
in which the test sees what the invariant cannot.

\subsection{Why no twist has to be inserted}\label{ssec:notwist}

That the two operators agree on the nose, rather than after a correction, has a
structural reason, and stating it is what makes the agreement something other
than a coincidence (\textbf{PROVEN}, \texttt{S17-05}).

System~$\mathsf C$'s leading coefficient factors as
\[
  3x^2+x+1=3(x-s_+)(x-s_-)=3\delta^2 z(z-1)=3\delta^2 D
\]
under $x=s_++\delta z$, with $D(z)=z(z-1)$ as in \S\ref{ssec:family}. Hence
System~$\mathsf C$'s own Wronskian coefficient is
\[
  P_{\mathsf C}=\frac{\delta\,(6x+1)}{3x^2+x+1}=\frac{D'}{D},
\]
which is \eqref{eq:PC}. The family's is $P=D'/D-b'/b$, and these coincide
exactly when $b$ is constant --- that is, when $b_1=0$, the condition
the match of \S\ref{ssec:located} already forced, and which
\S\ref{ssec:exhaustive} shows was implied by the coefficient equations in any
case. The two objects carry the same rank-$1$ twist, so there is no twist to
insert between them.

\subsection{Which scalar reduction, and why the choice is not
free}\label{ssec:route}

The comparison of \S\ref{ssec:pq} is between a scalar operator and a rank-$2$
system, so it presupposes an answer to a question that was for some time
contested: \emph{which} scalar equation the system is to be compared against.

\paragraph{There are two (\texttt{S17-09}, \textbf{PROVEN}).}
A rank-$2$ first-order system $y'=Ny$ admits two scalar reductions, the
equations satisfied by its two components. Derived generically by elimination,
with no formula assumed, route~$1$ (for $y_1$) has
$P=-\operatorname{tr}N-n_{12}'/n_{12}$ and route~$2$ (for $y_2$) has
$P=-\operatorname{tr}N-n_{21}'/n_{21}$. Fixing the orientation once, so that a
later reader cannot re-derive it backwards:
\begin{equation}\label{eq:route-diff}
  P_{\text{route }1}-P_{\text{route }2}
    =\frac{d}{dz}\log\frac{c}{b}
    =\frac{c'}{c}-\frac{b'}{b}.
\end{equation}
The proviso attached to \eqref{eq:route-diff} matters: in general the
difference is $\tfrac{d}{dz}\log(n_{21}/n_{12})$, and it equals
$\tfrac{d}{dz}\log(c/b)$ only because the convention $N=-A/D$ of
\S\ref{ssec:family} divides both off-diagonal entries by the same $D$. Both
routes are correct; neither is ``the'' scalar equation of the connection.

\paragraph{Route~$1$ is the one (\texttt{S17-10}, \textbf{PROVEN}).}
At the located point, route~$1$ reproduces System~$\mathsf C$ exactly --- $P-
P_{\mathsf C}=0$ and $Q-Q_{\mathsf C}=0$ symbolically --- while route~$2$
misses it by exactly $-c'/c=-1/(z-q)$. Both $b'=0$ and $q\notin\{0,1\}$ are
re-confirmed at the point.

\paragraph{A control independent of the derivation (\texttt{S17-11},
\textbf{VERIFIED}).}
Rather than compare formulas, one may solve the \emph{system} at an ordinary
point and ask which scalar equation the resulting $y_1$ satisfies. At
$z_0=\tfrac{3}{10}+\tfrac{i}{7}$, over two independent seeds, the route-$1$
residuals are $4.8\times10^{-61}$ and $6.9\times10^{-62}$ against route-$2$
residuals of $9.0$ and $4.0$. The control fires on both seeds, so the test can
in fact distinguish the routes.

\paragraph{Confirmation from a separately commissioned computation
(\texttt{S17-20}--\texttt{S17-23}).}
Because the choice of route had been contested, it was re-decided by a
computation isolated from the sessions that first made it. That computation
re-derived the two reductions from a generic $y'=Ny$ and obtained
\eqref{eq:route-diff} exactly, supplying the orientation and the denominator
proviso that the earlier derivation had left implicit (\texttt{S17-20},
\textbf{PROVEN}); neither of its two controls fired. At the point it was given,
the $b$-route reproduced System~$\mathsf C$ exactly over
$\mathbb Q(i,\sqrt3,\sqrt{11})$ --- by exact cancellation, never by comparing
radicals for equality --- with numeric residuals $2.0\times10^{-177}$ and
$2.5\times10^{-175}$ at $40$ digits over four sample points, while the
$c$-route missed by $-1/(z-q)$, numerically $1.45$, with the $Q$ difference
numerically $8.1$; the same test rejected a deliberately wrong point
$q\mapsto q+1$ (\texttt{S17-21}, \textbf{PROVEN}).

That computation also explains the earlier contested result rather than merely
overruling it (\texttt{S17-22}, \textbf{PROVEN}). The $c$-route residual
$-1/(z-q)$ is the simple zero of $c(z)=-(z-q)/(q-1)$ of \eqref{eq:fam-c} at
$z=q$, which is what \eqref{eq:route-diff} predicts when $b$ is constant at the
point. Since $-1/(z-q)$ has no zero in $z$, a comparison made against the
$c$-route returns a mismatch at every sample point and at every parameter
value: the earlier test could not have returned a match under any
circumstances. The residual was reached without the computation being told to
expect it.

Finally, the local data at infinity were checked independently
(\texttt{S17-23}, \textbf{PROVEN}). The leading matrix at infinity is
$\operatorname{diag}(t/2,-t/2)$ with distinct entries, so the $+t/2$ branch has
eigenvector $e_1+O(1/z)$ and its first component carries exponent
$\theta_\infty/2$ with no shift, for any $\theta_\infty$; System~$\mathsf C$
demands $-1-\sqrt3/18$, forcing $\theta_\infty=-2-\sqrt3/9$ with residual $0$
at $40$ digits, and the subdominant branch was predicted from
$\operatorname{ord}_\infty(n_{12})=2$ and matched exactly. Its scope is stated
in the claim itself and is repeated here: since $P$ and $Q$ already agree
identically, agreement of every local datum \emph{follows}. This is
corroboration by a second route, not an independent constraint, and it must not
be counted twice.

\begin{remark}[Two method defects, self-caught; \texttt{S17-25},
\textbf{PROVEN}]\label{rem:s20b-defects}
Recorded as findings about method, and deliberately not adopted as claims about
the mathematics. First, keeping $\theta_\infty$ symbolic inside the scalar
reduction and reading heads at infinity is unsound, because $\deg_z b$ drops
from $1$ to $0$ at exactly the value of $\theta_\infty$ of interest, so the
heads are discontinuous there. Second, the difference of the two exponents at
infinity is not gauge-invariant between the system and a scalar reduction, the
branches carrying powers $0$ and $-2$; this second defect returned the verdict
\emph{inconsistent} together with the value $\theta_\infty=\sqrt3/9$, wrong by
exactly the shift~$2$ of \S\ref{ssec:dictionary}. It was caught by an
arithmetic control --- the scalar exponents must sum to $-P_1$, and they summed
to $-1$ against System~$\mathsf C$'s $-2$. This is the converse of
Remark~\ref{rem:void}: a control that fired, on a real defect, inside a check
that would otherwise have returned a confident wrong answer.
\end{remark}

\subsection{The caveat, withdrawn}\label{ssec:withdrawn}

The ``up to rank-$1$ twist'' qualifier of \S\ref{ssec:membership-verdict} is
withdrawn (\texttt{S17-06}, \textbf{STRUCTURAL}). The defensible statement is
\begin{quote}
  the differential module of System~$\mathsf C$ is the module of the van der
  Put--Saito family at the located point of $U_0$,
\end{quote}
with no twist qualifier.

It is worth being exact about \emph{why} it is withdrawn, because the
conclusion was first reached for a reason that does not support it. Uniqueness
does not follow from trace-freeness; \S\ref{ssec:tracefree} refutes that. It
follows here from the direct verification of \S\ref{ssec:pq} that the full
operators agree. A correct conclusion does not certify the reasoning that
reached it, and each statement adopted in this paper is verified on its own
account.

The adjudication that settled the contested step is itself recorded with its
limits (\texttt{S17-24}, \textbf{STRUCTURAL}). The separately commissioned
computation was independent of the two sessions it adjudicated, but not of the
input file it was handed, which this programme wrote; and it was given the
point rather than locating one, so it is \emph{not} a second independent
location of a point of $U_0$. It tested the disputed step and nothing else,
which is what it was asked to do.

One bookkeeping note, recorded so that the citation trail is readable.
\texttt{S17-09} was amended in place to carry the orientation and the
denominator proviso of \eqref{eq:route-diff}; its label and its mathematics are
unchanged, and the prior text is preserved verbatim in the record
(\texttt{S17-26}, \textbf{STRUCTURAL}).

\medskip
\noindent
What \S\ref{sec:point} and this section together establish is an isomorphism of
differential modules at an explicitly located, provably unique pair of points,
and therefore a unique Painlev\'e~V leaf through System~$\mathsf C$. What they
do not establish is that the normalised representative of a projective class is
unique in general --- that remains \S\ref{sec:open}'s question --- and the
monodromy and Stokes data, which \S\ref{sec:monodromy} computes, are not used
anywhere above.
